% Part: computability
% Chapter: computability-theory
% Section: introduction

\documentclass[../../../include/open-logic-section]{subfiles}

\begin{document}

\olfileid{cmp}{thy}{int}
\olsection{Pambuka}

Cabang logika kang diarani \emph{teori komputabilitas}
nyinaoni perkara komputabilitas, utawa komputabilitas
relatif, saka fungsi lan himpunan. Pangaribawane
Kleene katon saka kasunyatan yen bidang iki biyen
diarani \emph{teori rekursi}, lan saiki jeneng loro-lorone
lumrah digunakake.

Fungsi~$f\colon \Nat \pto \Nat$ kita arani
\emph{komputabel parsial} yen bisa diitung ing
sawenehing model komputasi. Yen $f$ total, kita
cukup nyebut $f$~\emph{komputabel}. Relasi~$R$
kang fungsi karakteristike~$\Char{R}$ komputabel
uga diarani komputabel. Yen $f$ lan~$g$ fungsi
parsial, kita nulis $f(x) \fdefined$ kanggo
nyatakake yen $f$ ditegesi ing~$x$, yaiku
$x$~kalebu domain~$f$; lan $f(x) \fundefined$
kanggo nyatakake kosok baline, yaiku $f$
ora ditegesi ing~$x$. Kita nggunakake
$f(x) \simeq g(x)$ kanggo nyatakake yen
$f(x)$ lan $g(x)$ kalorone ora ditegesi,
utawa kalorone ditegesi lan padha.

Teori komputabilitas bisa ditliti tanpa kudu
ngrujuk marang model komputasi tartamtu.
Kanggo iki, dibuktekake yen ana fungsi
komputabel parsial universal~$\fn{Un}(k, x)$.
Iki ndadekake fungsi komputabel parsial bisa
dienumerasi. Kita nggunakake notasi~$\cfind{k}$
kanggo fungsi komputabel parsial unari kang
kaparingi urutan kaping~$k$, kang ditegesi
dening $\cfind{k}(x) \simeq \fn{Un}(k, x)$.
(Kleene nggunakake $\{ k \}$ kanggo tujuan
iki, nanging notasi kasebut saiki luwih
arang digunakake.) Kanthi rada luwih umum,
kita bisa ngenumerasi kanthi seragam fungsi
komputabel parsial kanthi aritas sembarang,
lan nggunakake $\cfind{k}[n]$ kanggo fungsi
rekursif parsial kaping~$k$ kanthi aritas~$n$.

Yen $f(\vec x, y)$ fungsi total utawa parsial,
mula $\umin{y}{f
(\vec x, y)}$ iku fungsi saka~$\vec x$ kang
ngasilake~$y$ paling cilik kang ndadekake
$f(\vec x, y) = 0$, kanthi syarat kabeh
$f(\vec x, 0)$, \dots, $f(\vec x, y-1)$
ditegesi; yen ora ana $y$ kaya mangkono,
$\umin{y}{f (\vec x, y)}$ ora ditegesi.
Yen $R(\vec x, y)$ relasi, $\umin{y}{R(\vec x, y)}$
ditegesi minangka~$y$ paling cilik kang
ndadekake $R(\vec x, y)$ bener; kanthi
tembung liya, $y$ paling cilik kang
ndadekake $1 \tsub \Char{R}(\vec x, y) = 0$.

Kanggo nuduhake yen fungsi iku komputabel,
ana rong cara kang bisa ditempuh:
\begin{enumerate}
\item Kanthi ketat: jlentrehna mesin Turing
  utawa fungsi rekursif parsial kanthi eksplisit,
  banjur tuduhna yen piranti iku ngitung fungsi
  kang dimaksud;
\item Kanthi informal: jlentrehna algoritma kang
  ngitung fungsi iku, banjur gunakna tesis Church.
\end{enumerate}
Ora ana wates kang cetha antarane rong cara iki;
jlentrehan algoritma kang rinci kudune menehi
katrangan cukup supaya lumayan cetha carane,
saora-orane sacara prinsip, ngrancang mesin
Turing utawa runtutan definisi rekursif parsial
kang trep. Definisi kang ketat kanthi lengkap
biasane ora akeh menehi pangerten, mula kita bakal
ngupaya nemokake imbangan antarane rong cara
iki; cekake, kita bakal ngupaya bukti kang
intuitif nanging tetep ketat,
kang nuduhake yen definisi kang presisi bisa disusun.

\end{document}
